%!TEX TS-program = xelatex
\documentclass{beamer}
\usepackage[french]{babel}
\usepackage[utf8]{inputenc}
\usepackage[T1]{fontenc}

\graphicspath{{figures/}} % Emplacement des figures.

\def\illustration{1} % Illustration en première page (1=activé, 0=désactivé).
					 % Penser à placer un fichier "illustration" dans le dossier des figures. Format préféré : 12,8cm par 4,8cm.
					 
\def\logo{1} % Logo en première page (1=activé, 0=désactivé).
			 % Penser à placer un fichier "logo" dans le dossier des figures.
						   
\def\outline{1} % Plan à chaque début de section (1=activé, 0=désactivé).

\usetheme{obsidienne} % Fichiers du thème obsidienne

\title{Obsidienne} % Titre
\subtitle{Un thème pour Beamer} % Sous-titre
\date[\today]{\today} % Date
\author[Th. G]{Thibault Giauffret\\ \href{http://thibault.giauffret.free.fr}{\texttt{thibault.giauffret.free.fr}}} % Auteur
\institute{ENS de Lyon} % Etablissement

\begin{document}
	
\begin{frame}[plain]
	\titlepage
\end{frame}

\section{Une section}
\begin{frame}
	\frametitle{Une première frame}
	\begin{block}{Bloc standard}
		Un bloc tout simple, par défaut.
		Un texte sur un fond de couleur qui 
		dépend, bien sûr, du thème choisi. 
	\end{block}

	\begin{alertblock}{Bloc standard}
		Un bloc tout simple, par défaut.
		Un texte sur un fond de couleur qui 
		dépend, bien sûr, du thème choisi. 
	\end{alertblock}

	\begin{exampleblock}{Bloc standard}
		Un bloc tout simple, par défaut.
		Un texte sur un fond de couleur qui 
		dépend, bien sûr, du thème choisi. 
	\end{exampleblock}
\end{frame}

\begin{frame}
\frametitle{Une seconde frame}
\begin{figure}[h!]
	\includegraphics[width=0.4\textwidth]{tiger}
	\caption{Une jolie image vectorielle}
\end{figure}
\end{frame}

\section{Une seconde section}
\begin{frame}
\frametitle{Une première frame}
\end{frame}

\begin{frame}
\frametitle{Une seconde frame}
\end{frame}

\section{Une seconde section}
\section{Une troisième section}
\section{Une quatrième section}
\section{Une cinquième section}

\end{document}